\documentclass[10pt,a4paper]{amsart}
\usepackage[margin=19mm]{geometry}
\usepackage{amsmath,amssymb,mathtools}
\usepackage[scr=rsfs,cal=euler]{mathalfa}
\usepackage{xfrac}
\usepackage[unicode,hidelinks,bookmarks=false]{hyperref}
\newcommand{\ie}{i.e.\ }
\newcommand{\cf}{cf.\ }
\newcommand{\resp}{respectively\ }
\newcommand{\Subsection}{\subsection}
\providecommand{\nobreakdash}{\nobreak\hbox{-}}
\setlength{\emergencystretch}{2em}
\multlinegap0pt
\usepackage{amssymb}
\usepackage{enumerate}
\usepackage[shortlabels]{enumitem}
\newtheorem{prop}{Proposition}
\DeclareMathOperator{\aut}{Aut} 
\DeclareMathOperator{\inn}{Inn} 
\DeclareMathOperator{\perm}{Sym}
\title{Characterizing finite solvable groups through the nilpotency probability}
\author{Andrea Lucchini}
\date{Source: arXiv:2604.04534v1 (CC BY 4.0)}
\begin{document}
\maketitle

\noindent\emph{Source and licence.} Characterizing finite solvable groups through the nilpotency probability. Version arXiv:2604.04534v1 (CC BY 4.0), distributed by arXiv under the Creative Commons Attribution 4.0 International licence, http://creativecommons.org/licenses/by/4.0/, as recorded in the arXiv licence metadata for this exact version and shown on the abstract page https://arxiv.org/abs/2604.04534v1. Full licence notice: \texttt{licenses/arxiv-2604.04534-CC-BY-4.0.txt}.\par\smallskip

\noindent\textit{Editorial scope.} Counted unit: the article's prop mono (source0.tex lines 192--196). The article's complete original proof of it is delivered with it (source0.tex lines 197--220, one \texttt{proof} environment of 24 lines). No mathematics is written, repaired or re-derived: the statement and the proof are the article's own lines, in the article's own notation, and no retained line is substituted or re-flowed.  No further source-local context is needed: every cross-reference of every retained line resolves inside the two delivered spans.  Nothing was omitted from any retained span, so the package carries no omission record.  No retained line contains a citation, so the wrapper carries no bibliography and no bibliography entry.  No retained line contains a figure environment, an image-inclusion command or drawing code, so no image is delivered and the package declares no asset dependency.  Editorial formatting only, in addition: an AMS \texttt{amsart} preamble that restates the article's own declarations and macros used by the retained lines, namely the environments \texttt{prop} on separate counters, together with the standard packages those lines need.  The retained environments print in this document as Proposition 1 in order of appearance, whereas the article prints the same items with the numbering of its own full document.  The complete native source of the article as distributed in the arXiv e-print for this exact version is bundled unchanged as \texttt{sources/197888/source0.tex}.
\begin{prop}\label{mono}
	Assume that a finite group $G$ contains a unique minimal normal subgroup, say $N$, and that $N$ is nonabelian. Then
	$$\nu(G)\leq \tilde\nu(S),$$
	where $S$ is a composition factor of $N.$
\end{prop}

\begin{proof}
Assume that $N=S^u,$ with $u\in \mathbb N.$ 	We may identify $G$ with a subgroup of the wreath product $\aut S\wr \perm(u).$  In this identification, any element of $G$ can be written in the form $(\alpha_1,\dots,\alpha_u)\rho,$ with $\alpha_1,\dots,\alpha_u \in \aut S$ and $\rho \in \perm(u).$ 
Moreover, we identify $S$ with the inner automorphism group $\inn(S)$ and therefore
$N=\{(s_1,\dots,s_u)\mid s_1,\dots,s_u \in S\}.$

Now let $$g_1=(a_1,\dots,a_u)\sigma,\quad g_2=(b_1,\dots,b_u)\tau \quad \in G.$$
Write $$\sigma=\sigma_1\cdots \sigma_c, \quad  \tau=\tau_1\cdots \tau_d,$$ as products of disjoint cyclic permutations and assume that 1 belongs to the supports of $\sigma_1$ and $\tau_1.$ Suppose
$$\sigma_1=(1,i_2,\dots,i_e), \quad \tau_1=(1,j_2,\dots,j_f).$$
Let $$X=\{(\alpha_1,\dots,\alpha_u)\rho \in G \mid \rho(1)=1\}.$$ The map $\phi\colon X\to \aut S$ sending $(\alpha_1,\dots,\alpha_u)\rho$ to $\alpha_1$ is a group homomorphism. Now let $n_1=(s_1,\dots,s_u), n_2=(t_1,\dots,t_u) \in N$. Then
$$(n_1g_1)^e, (n_2g_2)^f \in X$$ and
$$\phi((n_1g_1)^e)=s_1a_1s_{i_2}a_{i_2}\cdots s_{i_e}a_{i_e},\quad
\phi((n_2g_2)^f)=t_1b_1t_{j_2}b_{j_2}\cdots t_{j_f}b_{j_f}.
$$
Let
$$\gamma_1=a_{1}s_{i_2}a_{i_2}\cdots s_{i_e}a_{i_e},\quad \gamma_2=b_{1}t_{j_2}b_{j_2}\cdots t_{j_f}b_{j_f}.$$
If $(s_1,t_1)\notin  \mathcal N_{\gamma_1,\gamma_2}(\aut S,S),$ then
	$$\phi(\langle (n_1g_1)^{e}, (n_2g_2)^f\rangle)=
	\langle s_1\gamma_1,t_1\gamma_2\rangle$$ is not nilpotent. In particular $\langle n_1g_1, n_2g_2\rangle$ is not nilpotent. Thus, given $s_i, t_j$ for $i\neq 1$ and $j\neq 1$, there are at least $$|S|^2(1-\nu_{\gamma_1,\gamma_2}(\aut S, S))\geq |S|^2(1-\tilde\nu(S))$$ choices for 
	$s_1$ and $t_1$ such that $\langle n_1g_1, n_2g_2\rangle$ is not nilpotent. It follows that there
	are at least $|S|^{2u}(1-\tilde \nu(S))$ pairs $(n_1,n_2)\in N^2$ such that $\langle n_1g_1,n_2g_2\rangle$ is not nilpotent. This is equivalent to saying that $\nu_{g_1,g_2}(G,N)\leq \tilde \nu(S)$ for every $(g_1,g_2) \in G^2.$ But then
	$$\begin{aligned}\nu(G)|G|^2=&\sum_{(Ng_1,Ng_2)\in  (G/N)^2}\nu_{g_1,g_2}(G,N)\\&\leq \sum_{(Ng_1,Ng_2)\in (G/N)^2}|N|^2\tilde\nu(S)\\&\leq |G|^2\tilde\nu(S),
	\end{aligned}$$
	and consequently $\nu(G)\leq \tilde \nu(S).$
\end{proof}

\end{document}
